\section{USAJMO 2025}
        \begin{enumerate}
        	\item (\textit{USAJMO 2025}) Let $\mathbb Z$ be the set of integers, and let $f\colon \mathbb Z \to \mathbb Z$ be a function. Prove that there are infinitely many integers $c$ such that the function $g\colon \mathbb Z \to \mathbb Z$ defined by $g(x) = f(x) + cx$ is not bijective.
Note: A function $g\colon \mathbb Z \to \mathbb Z$ is bijective if for every integer $b$, there exists exactly one integer $a$ such that $g(a) = b$.


 

	\item (\textit{USAJMO 2025}) Let $k$ and $d$ be positive integers. Prove that there exists a positive integer $N$ such that for every odd integer $n>N$, the digits in the base-$2n$ representation of $n^k$ are all greater than $d$.


 

	\item (\textit{USAJMO 2025}) Let $m$ and $n$ be positive integers, and let $\mathcal R$ be a $2m\times 2n$ grid of unit squares.


 A domino is a $1\times2$ or $2\times1$ rectangle. A subset $S$ of grid squares in $\mathcal R$ is domino-tileable if dominoes can be placed to cover every square of $S$ exactly once with no domino extending outside of $S$. Note: The empty set is domino tileable.


 An up-right path is a path from the lower-left corner of $\mathcal R$ to the upper-right corner of $\mathcal R$ formed by exactly $2m+2n$ edges of the grid squares.


 Determine, with proof, in terms of $m$ and $n$, the number of up-right paths that divide $\mathcal R$ into two domino-tileable subsets.


 

	\item (\textit{USAJMO 2025}) Let $n$ be a positive integer, and let $a_0,\,a_1,\dots,\,a_n$ be nonnegative integers such that $a_0\ge a_1\ge \dots\ge a_n.$ Prove that

 \[\sum_{i=0}^n i\binom{a_i}{2}\le\frac{1}{2}\binom{a_0+a_1+\dots+a_n}{2}.\]
Note: $\binom{k}{2}=\frac{k(k-1)}{2}$ for all nonnegative integers $k$.


 

	\item (\textit{USAJMO 2025}) Let $H$ be the orthocenter of acute triangle $ABC$, let $F$ be the foot of the altitude from $C$ to $AB$, and let $P$ be the reflection of $H$ across $BC$. Suppose that the circumcircle of triangle $AFP$ intersects line $BC$ at two distinct points $X$ and $Y$. Prove that $C$ is the midpoint of $XY$.


 

	\item (\textit{USAJMO 2025}) Let $S$ be a set of integers with the following properties:


 $\bullet$ $\{ 1, 2, \dots, 2025 \} \subseteq S$.


 $\bullet$ If $a, b \in S$ and $\gcd(a, b) = 1$, then $ab \in S$.


 $\bullet$ If for some $s \in S$, $s + 1$ is composite, then all positive divisors of $s + 1$ are in $S$.


 Prove that $S$ contains all positive integers.


 

\end{enumerate}